Differential equations describe how fluids flow and how heat spreads, and the cost of solving one numerically grows with the resolution it needs: $N$ unknowns along each of $D$ axes make $N^D$ in all. A register of $n$ qubits holds a vector of $N=2^n$ amplitudes, and this exponential compression is the reason quantum algorithms for differential equations are proposed \cite{berry2014highorder,childsliu2020spectral,liu2021carleman}. What such an algorithm returns is a quantum state, from which measurement recovers only a few properties of the solution \cite{aaronson2015fineprint}, and what it costs on a fault-tolerant machine is the number of queries it makes to a circuit encoding the equation, multiplied by the gates in each query. Every route pays a ratio that the discretisation fixes before any circuit is built: the condition number of a linear system, the strength of a nonlinearity against the dissipation in a linearised equation, or the norm of an encoded Hamiltonian against its spectral gap. An accurate discretisation can still make that ratio large. For an equation in a fixed number of dimensions with a smooth solution, the most a quantum solver can gain over the best classical methods is polynomial \cite{montanaro2016fem}, so this thesis compares quantum formulations with one another and claims no advantage over classical solvers.

One recent route makes the solution a ground state. It descends from the differentiable quantum circuits of Kyriienko, Paine and Elfving, which train a parametrised circuit on the residual of the equation \cite{kyriienko2021differentiable}, and from the Chebyshev latent space of Paine, Elfving and Kyriienko, in which the solution's Chebyshev coefficients are the amplitudes of a state and differentiation acts on them as a fixed matrix \cite{paine2023physicsinformed}. Wu, Paine, Philip, Gentile and Kyriienko removed the training \cite{wu2025pihm}. In their \gls{pihm}, the equation and each of its data conditions become operators that annihilate the state exactly when it holds the solution's coefficients, and the sum of their Gram matrices is a Hamiltonian $H$ whose ground state is the solution. A spectral filter prepares that state directly, with nothing to optimise. Wu et al.\ show that every term of $H$ has a circuit of $\mathcal{O}(\mathrm{poly}(n))$ gates, so one application of the Hamiltonian is cheap.

A cheap application says nothing about how many applications a preparation needs. That number is the degree of the filter, and for a block encoding of $H$ at subnormalisation $\alpha_H\ge\|H\|$, a gap $\Delta$ above the ground state and a suppression $\varepsilon$ of the excited states, it grows as $d=\Theta\bigl(\sqrt{\alpha_H/\Delta}\,\log(1/\varepsilon)\bigr)$ (\Cref{eq:degree}). Wu et al.'s Discussion leaves open two problems that bear on it. They write that ``studying the spectral gap dependence is an important question for the future work'', and they suggest that the choice of basis be guided by ``the cost of differentiation'' \cite{wu2025pihm}, which they measure as circuit depth. The gap is the degree's denominator. On our reading the cost of differentiation sits in its numerator as well: each derivative eliminated through the Chebyshev differentiation matrix, which is dense and whose norm grows as $N^2$, multiplies the residual's norm by a factor of that order and the Hamiltonian's by its square, and the filter pays for that norm at every query (\Cref{sec:pihm-assessment}). We therefore ask how the number of queries grows with the resolution, and how much of that growth belongs to the encoder and how much to the Hamiltonian; whether that number can be lowered by changing the residual rather than the circuit, while keeping the solution; and whether the answer survives the step to nonlinear equations, where Wu et al.'s construction selects a state from a degenerate kernel by projection instead of preparing a unique one.

The two costs can be told apart on a single problem. Take the critically damped oscillator $f''+4f'+4f=0$, one of Wu et al.'s examples and here called R2a, at $n=8$ qubits, with the minimax edge filter, the polynomial of least degree that suppresses the excited states to $\varepsilon=\res[1,sci]{stateprep/R2a/standard/n8/T3/filter=minimax}{epsilon}$. With Wu et al.'s encoding of the differentiation matrix the filter's degree is $\res[2,sci]{stateprep/R2a/standard/n8/T3/filter=minimax/encoder=published}{degree}$. An exact encoding of the same Hamiltonian brings it down to $\res[2,sci]{stateprep/R2a/standard/n8/T3/filter=minimax}{degree}$, and no encoding of that Hamiltonian can do better than $\res[2,sci]{stateprep/R2a/standard/n8/T3/filter=minimax}{ideal_degree}$, the degree at the ideal-encoding bound $\alpha=\|R\|$, where $R$ is the residual operator whose Gram matrix is $H$. The fall from the first number to the second is the encoder's cost, and an exact encoding removes it. What remains is the Hamiltonian's cost, which only a different Hamiltonian with the same ground state can remove. The descriptor form developed in this thesis is such a Hamiltonian, and on the same problem its filter needs a degree of $\res{stateprep/R2a/descriptor/n8/T3/filter=minimax}{degree}$ as built and $\res{stateprep/R2a/descriptor/n8/T3/filter=minimax}{ideal_degree}$ at its own ideal bound. All five are resource estimates, each computed from a subnormalisation and the classically computed gap; the standard form's degrees lie far beyond the simulation budget. R2a is a regular problem, posed on one axis with polynomial coefficients, a leading coefficient without zeros on $[-1,1]$, analytic data and a solution unique up to scale, and it is where the reformulation gains most. On singular problems and on two axes it gains less, and in two cases the comparison reverses (\Cref{sec:contributions}).

\subsection{Contributions and Scope}
\label{sec:contributions}

The principal contribution is a reformulation of the residual, which we call the descriptor form. Chebyshev differentiation is dense, but it factorises into two banded matrices, $\G=\Bt^{-1}\Dt$, which are, up to a factor of two, the conversion and differentiation operators of Olver and Townsend's ultraspherical spectral method \cite{olver2013}. The descriptor form keeps the two factors and never forms $\G$. An equation of second or higher order carries its intermediate derivatives as unknowns, each tied to the next by a banded recurrence, in what we call the chain; a first-order equation or system is multiplied through by $\Bt$, in the mass form. In both layouts the residual's kernel is the standard form's, so the state prepared is the same solution. Carrying derivatives as unknowns so that high order does not compound the conditioning is the idea of \gls{fosls} \cite{cai1994fosls}. What is new here is the transfer of these two classical ideas to a block-encoded Hamiltonian, and a measurement of what it does to the cost of preparing the ground state. The Hamiltonian's norm falls from $\Theta(N^{4k})$ for an order-$k$ equation with constant coefficients to $\Theta(N^2)$, and the residual compiles to a single \gls{lcu} with $\mathcal{O}(\log n)$ ancilla qubits, $\lceil\log_2(n+4)\rceil+8$ for the oscillators' equation block. On a regular problem the gap is proved bounded below independently of $N$ (\Cref{thm:flat-gap}), and the degree grows as $\tilde\Theta(N)$, where for constant coefficients the standard form's grows as $\tilde\Theta(N^{2k})$ even under the best exact encoding.

The reformulation is not free. The derivatives need a block register, and at the sizes built a descriptor query costs more two-qubit gates than a standard one on constant-coefficient problems, so the advantage there lies in the number of queries alone. The carried derivatives also erode the gap wherever the equation fails to control them. On singular problems, whose leading coefficient vanishes somewhere on $[-1,1]$, an a priori block rescaling restores the gap without ancilla, lengthening each query by a factor that the equation fixes, and the degree then grows as about $N^2$; on Legendre's equation the standard form needs the smaller degree at the ideal bound. On two axes the rescaled gap still falls with $N$, and on the heat equation at four qubits per axis the standard form needs the lower degree. A double zero of the leading coefficient is beyond any rescaling. Carrying the intermediate derivatives instead of eliminating them therefore makes the residual banded and keeps its kernel, lowering the Hamiltonian's norm from $\Theta(N^8)$ to $\Theta(N^2)$ for a second-order equation; the price is a block register and a gap that the derivative blocks erode, which an a priori rescaling restores on one axis without ancilla (\Cref{sec:descriptor,sec:results-scaling,sec:results-gap}).

The reformulation rests on four further pieces of work, listed in the order the thesis presents them.
\begin{enumerate}
    \item We state Wu et al.'s construction once, as the standard form, in which the state holds the field alone and every derivative is eliminated through $\G$. Where a printed choice would make a problem inaccurate, answer-dependent or ill posed, the standard form replaces it at the point where the construction meets it (\Cref{tab:deviations}); the replacements change the data and leave $\G$ and the Gram structure alone, so the baseline keeps Wu et al.'s cost. With the printed choices restored, the construction reproduces the scale factors of all eleven of their \gls{ode} examples to within one unit of the last printed digit. Their partial differential and nonlinear examples do not reproduce, and the hypotheses tested are recorded (\Cref{sec:standard}, \Cref{app:reproduction}).
    \item An exact factorisation of the differentiation matrix gives an encoding of $\G$ with $\mathcal{O}(n^2)$ gates and $n+4$ ancilla qubits, within a factor $\res[3]{derivative/n10}{structured_ratio}$ of its norm at $n=10$, against $\res[2,sci]{derivative/n10}{published_ratio}$ for Wu et al.'s approximate encoding. It exists to make the comparison fair: for constant coefficients it holds the encoder's cost to a constant, so whatever growth remains in the standard form's degree belongs to its Hamiltonian. For a lifted variable coefficient it is loose, and there only the ideal bound separates the two costs (\Cref{sec:standard-factorisation,sec:results-scaling}).
    \item Both forms are carried beyond the single real field that Wu et al.\ treat. The fields of a coupled first-order system share one state, a block each, with vector data entering as rows that fix a direction and leave the scale free, and a coupling compiled as a sum of matchings; the Fourier bases for which Wu et al.\ set complex-valued states aside are carried on real ones. For nonlinear equations we show that Wu et al.'s doubled space has a kernel of dimension at least $2^{2n-1}$, so a filter only projects onto it, while a Carleman lift, itself a coupled linear system, has a one-dimensional kernel at every order run and is prepared in both forms, up to two-dimensional incompressible Navier--Stokes on $12$ to $48$ Fourier modes by exact emulation and resource estimate. A trusted kernel, one-dimensional and well isolated, turns out not to certify that the lift has converged (\Cref{sec:extensions,sec:results-nonlinear,sec:results-ns}).
    \item One preparation pipeline serves both forms, so that a difference measured between them is a difference between their residuals. It block-encodes $H$ through a reflection built from the stacked residual and filters the quantum walk of that reflection by \gls{gqsp} with the minimax edge filter, which a study in both forms preferred to Wu et al.'s imaginary-time filter; its parameters are fixed a priori and its initial states chosen without knowledge of the answer. Every Hamiltonian of at most $24$ qubits is verified whole against the operator assembled from the problem, every lift and coupling term by term at any size, and wherever both run, the simulated circuit equals the exact emulation within $10^{-10}$ before normalising on every run of the minimax filter (\Cref{sec:pipeline,sec:results-faithful}).
\end{enumerate}

Every result in the thesis is a classical computation, and nothing ran on quantum hardware. The evidence comes in three tiers, named wherever a number is quoted: circuit simulation (T1), which reaches $\res{stateprep/R4d/descriptor/n3/T1/filter=minimax}{n_qubits}$ qubits; exact emulation (T2), provably equal to the circuit's post-selected output, which reaches filter degrees of $2^{20}$; and resource estimates (T3) beyond that, a circuit too wide to simulate being reported as beyond the simulation budget and never as out of the algorithm's reach. The field is decoded classically from the prepared state's amplitudes, which a quantum computer does not expose, and scaled by a known value of the solution; Wu et al.'s interferometric readout is built and verified, with a correction for a post-selection that succeeds with probability below one, only at small $n$ on R1, R2a and R5a, and costing it is further work. The edge of the excited spectrum, which the filter's degree needs, is taken from the classical reference solution. Neither form is polylogarithmic in $N$, and no advantage over classical solvers is claimed: on one axis a classical spectral solver needs only $\mathcal{O}(N)$ operations \cite{olver2013}. The floor on the standard form's degree is shown for polynomial filters of the walk, and since Lin and Tong's optimality result for ground-state preparation assumes access to $H$ rather than to its residual \cite{lintong2020groundstate}, whether it binds every algorithm is left open. No solution is said to be prepared from a degenerate formulation, no complex-valued field is prepared natively, and a trusted kernel is never read as a converged lift.

The work departs from its research proposal (\Cref{app:proposal}), which set out to carry Wu et al.'s method to complex-valued nonlinear equations, among them the nonlinear Schr\"odinger and complex Ginzburg--Landau equations, through a generalised \gls{qsvt}, and to add data-informed terms to its Hamiltonian. Reproducing and costing Wu et al.'s examples showed that the query count, whose growth with the resolution their Discussion leaves open, grows as a power of $N$, through the encoding of the dense differentiation matrix and through the norm that matrix gives the Hamiltonian. Any extension would inherit that cost, so the thesis settles it first. The filter is applied by \gls{gqsp}, which already lifts the two restrictions the proposal sought a generalised transform for, real coefficients and a definite parity, and a complex residual would need no further transform, since the filter acts only on the Hermitian $R^\dagger R$. Of the proposed equations only Burgers' is treated, joined by two-dimensional Navier--Stokes; the Schr\"odinger and Ginzburg--Landau equations and the data-informed Hamiltonians are left to further work.

\subsection{Outline}
\label{sec:outline}

\Cref{sec:background}, the background and literature review, follows the ratio a discretisation fixes through the quantum routes to differential equations, the spectral methods on which the construction rests and the quantum primitives that turn a ratio into a count of queries, judging each. \Cref{sec:pihm} presents the \gls{pihm} as Wu et al.\ published it, as prior art, assesses it, and ends on the three questions that the rest of the thesis answers. \Crefrange{sec:standard}{sec:pipeline} are the thesis's own model. \Cref{sec:standard}, the standard form and its cost, states the construction once and locates its cost; \Cref{sec:descriptor}, the descriptor reformulation, changes the residual and accounts for what the change costs; \Cref{sec:extensions}, on systems and nonlinearity, carries both forms from one real field to coupled and nonlinear ones; and \Cref{sec:pipeline}, on preparing and verifying the ground state, describes the pipeline through which both forms are tested and the design of the numerical experiment. \Cref{sec:results}, the numerical study, gives the evidence, one section per question. \Cref{sec:discussion} discusses the mechanism behind the findings and their limits, and \Cref{sec:conclusions}, the conclusions and future work, gives the verdicts and the work they leave open. The appendices hold the proofs and constructions (\Cref{app:cheb} to \Cref{app:verification}), the catalogue of benchmark problems (\Cref{app:catalogue}), the reproduction and audit of Wu et al.'s examples (\Cref{app:reproduction}), the software and records (\Cref{app:software}), and the research proposal with its deviations (\Cref{app:proposal}).
